% =====================================================================
%  Lecture 22 -- Quality of Service in 5G                 (English edition)
% =====================================================================
\documentclass[11pt,a4paper]{article}
\providecommand\ACNroot{../..}
\input{\ACNroot/templates/lecture-en}

\ACNset{number}{22}
\ACNset{title}{Quality of Service (QoS) in 5G}
\ACNset{shorttitle}{5G QoS}
\ACNset{titleother}{کیفیت خدمات (QoS) در نسل پنجم}
\ACNset{part}{Part IV -- 5G system and procedures (Lectures 14, 19--22)}
\ACNset{time}{6--8 hours (3 h reading, 3--5 h exercises)}
\ACNset{level}{Advanced}
\ACNset{prereq}{Lectures 7, 9, 13, 20 and 21; M/M/1 queue}
\ACNset{gen}{5}

\begin{document}
\ACNcover

% ---------------------------------------------------------------------
\section{Learning objectives}
After studying this lecture you should be able to:
\begin{objectives}
  \item \textbf{Explain} the 5G QoS model: QoS flows, QFI, QoS profiles, QoS rules and their mapping to radio bearers.
  \item \textbf{Interpret} standardised 5QI characteristics (resource type, priority, packet delay budget, packet error rate,
        maximum data burst volume).
  \item \textbf{Trace} how policy from the PCF becomes QoS flows enforced in UE, RAN and UPF.
  \item \textbf{Compute} delay budgets, HARQ needs for a packet-error target and the effect of priority scheduling.
  \item \textbf{Compare} the 5G flow-based model with LTE's bearer-based model.
\end{objectives}

% ---------------------------------------------------------------------
\section{Prerequisite knowledge}
You need LTE scheduling and HARQ (Lecture~7), the EPS bearer and QCI model (Lecture~9), PFCP rules (Lecture~13), PDU session
establishment (Lecture~20) and slicing (Lecture~21). The priority model extends the M/M/1 queue of Lecture~12.

\paragraph{Key terms} The vocabulary of this lecture with its Persian equivalents:
\begin{keyterms}
QoS flow & \fa{جریان کیفیت خدمات} & Finest granularity of QoS differentiation in a PDU session. \\
QFI & \fa{شناسهٔ جریان QoS} & 6-bit identifier of a QoS flow, carried in N3 encapsulation. \\
5QI & \fa{شاخص QoS نسل پنجم} & Pointer to a set of standardised QoS characteristics. \\
QoS profile & \fa{نمایهٔ QoS} & Parameters given to the RAN per QoS flow (5QI, ARP, GFBR, MFBR, \ldots). \\
QoS rule & \fa{قاعدهٔ QoS} & UE-side packet filters mapping uplink traffic to QoS flows. \\
GBR / non-GBR & \fa{با / بدون نرخ تضمین‌شده} & Resource types; delay-critical GBR for URLLC. \\
PDB / PER & \fa{بودجهٔ تأخیر / نرخ خطای بسته} & Upper bounds on delay and loss for a 5QI. \\
MDBV & \fa{بیشینهٔ حجم انفجار داده} & Largest burst to be served within the PDB (delay-critical GBR). \\
SDAP & \fa{پروتکل سازگاری دادهٔ خدمت} & RAN layer mapping QoS flows to data radio bearers. \\
Reflective QoS & \fa{QoS بازتابی} & UE derives uplink QoS rules from downlink packets. \\
\end{keyterms}

% ---------------------------------------------------------------------
\section{Introduction and motivation}
Different applications need different treatment: a voice call needs low delay and tolerates some loss; a file transfer needs no
loss but tolerates delay; a robot controller needs both low delay and extremely low loss; a background update needs neither. A
mobile network cannot satisfy all of these with a single best-effort pipe, because the scarce resource -- radio capacity -- must be
shared under load. QoS mechanisms decide which packets go first, which are guaranteed a rate, and which may be dropped.

LTE solved this with \emph{bearers}: one per QoS class, each requiring signalling to set up along the whole path. 5G replaces bearers in
the core with lightweight \emph{QoS flows}, marked per packet and mapped flexibly to radio bearers by the RAN \cite{ts23501}. The
resulting model is more scalable, supports URLLC with new delay-critical characteristics, and integrates with slicing (Lecture~21): a
slice chooses the network, QoS flows the treatment inside it.

% ---------------------------------------------------------------------
\section{Fundamental concepts}

\subsection{QoS flows}
A PDU session contains one or more QoS flows, each identified by a QFI (\cref{fig:model}). Every PDU session has a default QoS flow
(non-GBR) that carries traffic not matching any other rule. Classification happens at the edges:
\begin{itemize}
  \item \textbf{downlink}: the UPF classifies packets with PDRs (packet filters from PCC rules) and marks them with the QFI in the
        N3 encapsulation header;
  \item \textbf{uplink}: the UE classifies packets with \emph{QoS rules} (packet filters with precedence) and the SDAP layer marks them;
  \item \textbf{RAN}: the gNB maps QoS flows to data radio bearers (DRBs) as it sees fit -- several flows may share a DRB.
\end{itemize}
Because QFI marking is per packet, adding a QoS flow does not require a new tunnel: all flows of a PDU session share one N3 tunnel.

\begin{acnfigure}
  \resizebox{\linewidth}{!}{\acnfig{lecture-22/qos-model}}
  \caption{5G QoS model: classification in UE and UPF, QFI marking on N3, flow-to-DRB mapping by SDAP in the gNB.}
  \label{fig:model}
\end{acnfigure}

\subsection{QoS parameters}
\begin{definitionbox}{QoS profile and characteristics}
The SMF gives the RAN a \emph{QoS profile} per QoS flow: the 5QI, the allocation and retention priority (ARP), and for GBR flows the
guaranteed and maximum flow bit rates (GFBR, MFBR), optionally notification control and a maximum packet loss rate. The 5QI points to
\emph{QoS characteristics}: resource type (non-GBR, GBR, delay-critical GBR), priority level, packet delay budget (PDB), packet error rate
(PER), averaging window and, for delay-critical GBR, the maximum data burst volume (MDBV) \cite{ts23501}.
\end{definitionbox}
Aggregate limits apply to non-GBR traffic: the \emph{Session-AMBR} per PDU session (enforced by UPF and UE) and the \emph{UE-AMBR} across a
UE's sessions (enforced by the RAN).

\begin{table}[htbp]
\centering\small
\caption{Selected standardised 5QI values (TS 23.501, Table 5.7.4-1) \cite{ts23501}. Lower priority values mean higher priority.}\label{tab:5qi}
\rowcolors{2}{acnmist}{white}
\begin{tabularx}{\linewidth}{@{}c l c c c Y@{}}
\toprule
\rowcolor{white}5QI & Resource type & Priority & PDB & PER & Example services \\
\midrule
1  & GBR                & 20 & 100\,ms & $10^{-2}$ & Conversational voice \\
2  & GBR                & 40 & 150\,ms & $10^{-3}$ & Conversational video (live) \\
3  & GBR                & 30 & 50\,ms  & $10^{-3}$ & Real-time gaming, V2X messages \\
4  & GBR                & 50 & 300\,ms & $10^{-6}$ & Non-conversational video (buffered) \\
65 & GBR                & 7  & 75\,ms  & $10^{-2}$ & Mission-critical push-to-talk voice \\
5  & Non-GBR            & 10 & 100\,ms & $10^{-6}$ & IMS signalling \\
6  & Non-GBR            & 60 & 300\,ms & $10^{-6}$ & Buffered video, TCP applications \\
7  & Non-GBR            & 70 & 100\,ms & $10^{-3}$ & Voice, live video, interactive gaming \\
9  & Non-GBR            & 90 & 300\,ms & $10^{-6}$ & Default bearer traffic (web, e-mail) \\
82 & Delay-critical GBR & 19 & 10\,ms  & $10^{-4}$ & Discrete automation (MDBV 255 bytes) \\
85 & Delay-critical GBR & 21 & 5\,ms   & $10^{-5}$ & Electricity distribution, high voltage (MDBV 255 bytes) \\
\bottomrule
\end{tabularx}
\end{table}

\subsection{Packet delay budget and packet error rate}
The PDB is an upper bound on the delay between UE and the UPF's N6 termination, for a given confidence; it includes a core-network part
and an access-network part. The PER is an upper bound on the rate of packets processed by the link layer but not delivered to the upper
layer. For GBR flows, packets delayed beyond the PDB count as lost (for delay-critical GBR, they are counted in the PER), so PDB and PER
together define the reliability target of a URLLC flow.

% ---------------------------------------------------------------------
\section{Detailed technical explanation}

\subsection{From policy to QoS flows}
QoS starts with policy (\cref{fig:setup}). An application function (e.g.\ the IMS for a voice call) asks the PCF for treatment of a media
flow; the PCF produces PCC rules with 5QI, ARP and bit rates and sends them to the SMF. The SMF binds each PCC rule to a QoS flow -- an
existing one with the same parameters or a new one with a new QFI -- and installs: PDRs and QERs in the UPF (N4), the QoS profile in the gNB
(N2), and the QoS rule in the UE (NAS). The gNB performs admission control for GBR flows and configures DRBs and the SDAP mapping.

\begin{acnfigure}
  \resizebox{\linewidth}{!}{\acnfig{lecture-22/qos-setup}}
  \caption{Establishing a GBR QoS flow triggered by an application function (simplified).}
  \label{fig:setup}
\end{acnfigure}

\subsection{Reflective QoS}
For applications with many short flows (e.g.\ web browsing with different priorities) signalling a QoS rule for each would be costly. With
\emph{reflective QoS} the UPF sets a reflective QoS indicator (RQI) in downlink packets; the UE then creates an uplink QoS rule mirroring the
downlink packet's addresses and ports, with the same QFI, without NAS signalling.

\subsection{Enforcement in the RAN}
The gNB scheduler translates QoS profiles into scheduling decisions: priority levels order the queues, GFBR is guaranteed over the averaging
window, PDB determines how long a packet may wait (and which numerology, mini-slots or configured grants are needed), and PER determines the
target BLER and HARQ/repetition strategy. QoS notification control lets the gNB inform the SMF when a GFBR can no longer be guaranteed, so that
the application can adapt (e.g.\ lower video resolution).

\subsection{QoS monitoring}
Since Release~16 the network can measure packet delay per QoS flow between UE, gNB and UPF and report it to the SMF and, via the NEF, to
application functions -- the basis for SLAs on URLLC slices and for applications that adapt to measured latency.

% ---------------------------------------------------------------------
\section{Mathematical models}

\subsection{Splitting the delay budget}\label{sec:pdb}
\paragraph{Model} The PDB is shared between the core network (UPF to gNB, $D_{\mathrm{CN}}$) and the access network ($D_{\mathrm{AN}}$):
\begin{equation}\label{eq:pdb}
  D_{\mathrm{AN}} = \mathrm{PDB} - D_{\mathrm{CN}} .
\end{equation}
\paragraph{Example} 5QI~1 (100\,ms) with an assumed core part of 20\,ms leaves 80\,ms for scheduling, HARQ and transmission -- comfortable. 5QI~85
(5\,ms) with a UPF 50\,km away (0.25\,ms one way in fibre plus processing, say 1\,ms) leaves 4\,ms: the UPF must be at the edge and the radio
must use short slots.

\subsection{Rate needed for a data burst}\label{sec:mdbv}
\paragraph{Model} For a delay-critical GBR flow, a burst of up to MDBV bytes must be delivered within the access-network budget. The
instantaneous rate the scheduler must provide during the burst is
\begin{equation}\label{eq:mdbv}
  R_{\mathrm{burst}} \ge \frac{8\,\mathrm{MDBV}}{D_{\mathrm{AN}}} .
\end{equation}
\paragraph{Example} 5QI~82 (MDBV 255\,bytes, PDB 10\,ms) with $D_{\mathrm{AN}} = 9$\,ms: 227\,kbit/s; 5QI~85 (255\,bytes, PDB 5\,ms) with
$D_{\mathrm{AN}} = 4$\,ms: 510\,kbit/s. The rates are small; the challenge is delivering them \emph{reliably} within milliseconds.

\subsection{HARQ transmissions for a PER target}\label{sec:harq}
\paragraph{Model} If each transmission of a block fails independently with probability $p$ (the target BLER), the residual error after $n$
transmissions is $p^n$. To reach a target PER,
\begin{equation}\label{eq:harq}
  n \ge \frac{\ln \mathrm{PER}}{\ln p},
\end{equation}
and all $n$ transmissions must fit into $D_{\mathrm{AN}}$.
\paragraph{Example} PER $= 10^{-5}$: with $p = 0.1$, $n = 5$ transmissions (\cref{fig:harq}); with a conservative MCS giving $p = 0.01$, $n = 3$. If one
HARQ round trip takes 4 slots of 0.125\,ms (120\,kHz) = 0.5\,ms, five transmissions need 2--2.5\,ms -- feasible within 4\,ms; at 15\,kHz (1\,ms slots)
they would need about 20\,ms -- impossible. URLLC therefore combines robust MCS (low $p$), short slots and sometimes blind repetitions.
\paragraph{Limitation} Errors are not independent when the channel is in a deep fade; diversity (frequency, antennas, multi-connectivity) is
needed to make the independence assumption approximately valid.

\begin{acnfigure}
  \acnfig{lecture-22/harq-per}
  \caption{Residual packet error rate after $n$ transmissions (\cref{eq:harq}).}
  \label{fig:harq}
\end{acnfigure}

\subsection{Priority scheduling}\label{sec:prio}
\paragraph{Model} Two classes share a link modelled as M/M/1 with service rate $\mu$; class~1 (load $\rho_1$) has preemptive priority over
class~2; total load $\rho$. The mean sojourn times are
\begin{equation}\label{eq:prio}
  T_1 = \frac{1}{\mu(1-\rho_1)}, \qquad
  T_2 = \frac{1}{\mu(1-\rho_1)(1-\rho)},
\end{equation}
compared with $T = 1/(\mu(1-\rho))$ for FIFO.
\paragraph{Example} $\mu = 1000$\,packets/s, $\rho_1 = 0.2$: $T_1 = 1.25$\,ms regardless of the total load; at $\rho = 0.9$, $T_2 = 12.5$\,ms against 10\,ms for FIFO
(\cref{fig:prio}).
\paragraph{Interpretation} Priority protects the high-priority class completely from the load of lower classes, at a moderate cost to the others --
as long as the high-priority load itself stays small. This is why GBR and delay-critical flows need admission control: without it, too much
``high-priority'' traffic destroys the protection.

\begin{acnfigure}
  \acnfig{lecture-22/priority-delay}
  \caption{Mean delay with preemptive priority and with FIFO (\cref{eq:prio}, $\rho_1 = 0.2$, $\mu = 1000$\,s$^{-1}$).}
  \label{fig:prio}
\end{acnfigure}

% ---------------------------------------------------------------------
\section{Network architecture and diagrams}
QoS is enforced at three points: the UE (uplink classification and Session-AMBR), the gNB (scheduling per QoS flow and DRB, UE-AMBR, GFBR) and the
UPF (downlink classification, gating, MFBR and Session-AMBR, usage reporting). The PCF decides; the SMF translates decisions into rules for each
point (\cref{fig:model,fig:setup}). The transport network between gNB and UPF should map QFIs or 5QIs to DSCP classes so that it does not undo the
differentiation.

% ---------------------------------------------------------------------
\section{Step-by-step procedures}

\begin{procedurebox}{Adding a GBR QoS flow for a voice call}\label{proc:gbr}
\begin{steps}
  \item The IMS (AF) sends the call's media description to the PCF (Npcf\_PolicyAuthorization).
  \item The PCF derives a PCC rule (5QI~1, GFBR/MFBR from the codec, ARP) and notifies the SMF.
  \item The SMF allocates a QFI, installs a PDR and QER in the UPF and sends the QoS profile to the gNB and the QoS rule to the UE through the AMF.
  \item The gNB admits the flow, maps it to a DRB with suitable RLC mode and scheduling, and confirms; voice packets now receive GBR treatment.
  \item At call end the PCF removes the rule; the SMF releases the QoS flow in UPF, RAN and UE.
\end{steps}
\end{procedurebox}

\begin{procedurebox}{Checking a URLLC QoS design}\label{proc:urllc}
\begin{steps}
  \item From the 5QI, read PDB, PER and MDBV; subtract the core-network delay to obtain $D_{\mathrm{AN}}$ (\cref{eq:pdb}).
  \item Compute the burst rate from MDBV (\cref{eq:mdbv}) and the number of transmissions from PER and BLER (\cref{eq:harq}).
  \item Choose numerology, mini-slots or configured grants so that all transmissions fit into $D_{\mathrm{AN}}$.
  \item Check the priority load: the high-priority traffic must remain a small fraction of capacity (\cref{eq:prio}); configure admission control.
\end{steps}
\end{procedurebox}

% ---------------------------------------------------------------------
\section{Worked examples}

\begin{examplebox}{Choosing 5QIs}\label{ex:5qi}
\textbf{Problem.} Assign 5QIs to: VoNR voice, IMS SIP signalling, a live video call, web browsing, factory motion control.

\textbf{Solution.} Voice 1; SIP 5; live video 2 (or 7 if non-GBR); web 9 (default) or 6/8; motion control a delay-critical GBR value such as 82 (or 85 if
5\,ms is required), with the UPF at the edge.
\end{examplebox}

\begin{examplebox}{URLLC feasibility}\label{ex:urllc}
\textbf{Problem.} A flow with 5QI~85 carries 200-byte bursts. The UPF is on site (core delay 0.5\,ms). With 30\,kHz SCS (0.5\,ms slots), a HARQ round trip
of 4 slots and $p = 0.01$, can the PER be met?

\textbf{Solution.} $D_{\mathrm{AN}} = 4.5$\,ms; burst rate $8\times200/0.0045 = 356$\,kbit/s; $n = \ln10^{-5}/\ln0.01 = 2.5 \to 3$ transmissions needing $3\times2 = 6$\,ms -- too long.
Options: mini-slots or 120\,kHz SCS (HARQ RTT 0.5\,ms: 1.5\,ms total), or blind repetitions within one slot window.
\end{examplebox}

\begin{examplebox}{Priority protection}\label{ex:prio}
\textbf{Problem.} With \cref{eq:prio}, what happens to $T_1$ if the high-priority load grows from 0.2 to 0.6, and to $T_2$ at total load 0.9?

\textbf{Solution.} $T_1 = 1/(1000\times0.4) = 2.5$\,ms (doubled); $T_2 = 1/(1000\times0.4\times0.1) = 25$\,ms. High-priority traffic must itself be limited.
\end{examplebox}

% ---------------------------------------------------------------------
\section{Technology comparison}
\begin{table}[htbp]
\centering\small
\caption{LTE bearer-based QoS versus 5G flow-based QoS.}\label{tab:cmp}
\rowcolors{2}{acnmist}{white}
\begin{tabularx}{\linewidth}{@{}l Y Y@{}}
\toprule
\rowcolor{white} & LTE (EPS bearers) & 5G (QoS flows) \\
\midrule
Unit of QoS          & EPS bearer (radio + S1 + S5 bearer) & QoS flow, marked by QFI \\
Core tunnels         & One per bearer                      & One per PDU session \\
Radio mapping        & One radio bearer per EPS bearer     & Flexible flow-to-DRB mapping (SDAP) \\
QoS index            & QCI (1--9, 65--80)                  & 5QI, including delay-critical GBR \\
Uplink classification & UL TFT, signalled                  & QoS rules, signalled or reflective \\
New flow             & Dedicated bearer set-up end to end  & QoS rule/profile update, no new tunnel \\
Delay monitoring     & Limited                             & Per-flow QoS monitoring (Rel-16) \\
\bottomrule
\end{tabularx}
\end{table}

\paragraph{Same ideas, lighter mechanics} The QoS parameters (priority, PDB, PER, ARP, GBR/MBR, AMBR) are almost identical to LTE's; what changed is the
mechanism -- marking instead of tunnels, flexible RAN mapping, and reflective rules -- which makes fine-grained QoS affordable.

% ---------------------------------------------------------------------
\section{Key takeaways}
\begin{takeawaybox}
\begin{itemize}
  \item 5G QoS is flow-based: QFI marking per packet, one N3 tunnel per PDU session, SDAP maps flows to DRBs.
  \item 5QI gives resource type, priority, PDB, PER (and MDBV for delay-critical GBR); ARP and GFBR/MFBR complete the QoS profile.
  \item PCF $\to$ SMF $\to$ (UPF via N4, gNB via N2, UE via NAS): policy becomes enforceable rules.
  \item $D_{\mathrm{AN}} = \mathrm{PDB} - D_{\mathrm{CN}}$; burst rate $\ge 8\,\mathrm{MDBV}/D_{\mathrm{AN}}$; transmissions $n \ge \ln\mathrm{PER}/\ln p$.
  \item Priority protects high-priority traffic only if its own load is limited: $T_1 = 1/(\mu(1-\rho_1))$.
\end{itemize}
\end{takeawaybox}

% ---------------------------------------------------------------------
\section{Exercises}

\begin{exercisebox}{1}{Flows and bearers}\label{exr:flows}
A UE has an Internet PDU session with three QoS flows. How many N3 tunnels and how many DRBs does it need at least and at most?
\end{exercisebox}

\begin{exercisebox}{1}{Reflective QoS}\label{exr:rq}
Explain why reflective QoS saves signalling and give one situation where it should not be used.
\end{exercisebox}

\begin{exercisebox}{2}{Budget split}\label{exr:pdb}
For 5QI~3 (PDB 50\,ms) and 5QI~82 (PDB 10\,ms), compute $D_{\mathrm{AN}}$ when the UPF is 600\,km away (5\,\textmu s/km, plus 2\,ms processing) and when it is local
(0.5\,ms). Which 5QI forces an edge UPF?
\end{exercisebox}

\begin{exercisebox}{2}{HARQ}\label{exr:harq}
A flow needs PER $10^{-6}$. Compute the number of transmissions for $p = 0.1$ and $p = 0.001$, and the time needed with a HARQ round trip of 1\,ms. What does this
imply for the choice of MCS?
\end{exercisebox}

\begin{exercisebox}{3}{Admission control}\label{exr:adm}
A cell carries 5QI~1 voice (GBR 40\,kbit/s per call at the radio, assumed to cost 0.5\,\% of cell capacity per call). Using \cref{eq:prio} with voice as class 1,
how many calls can be admitted so that non-GBR traffic at total load 0.9 sees a mean delay at most twice its FIFO value? Interpret.
\end{exercisebox}

\begin{exercisebox}{3}{End-to-end QoS}\label{exr:e2e}
A cloud-gaming application needs 30\,ms RTT and 50\,Mbit/s downlink. Propose 5QI(s), UPF placement, slice and transport measures, and explain which nodes enforce
which parameter.
\end{exercisebox}

\subsection*{Solutions}
\addcontentsline{toc}{subsection}{Solutions}

\begin{solutionbox}{\ref{exr:flows}}
One N3 tunnel (per PDU session) in all cases. DRBs: at least one (all three flows mapped to one DRB) and at most three (one per flow); the gNB decides.
\end{solutionbox}

\begin{solutionbox}{\ref{exr:rq}}
The UE derives uplink rules from downlink packets marked with the RQI, so no NAS signalling is needed per flow -- valuable for many short-lived flows. It
should not be used where the uplink treatment must differ from what the downlink packet implies, or for security-sensitive policy where the network must
control exactly which uplink packets get privileged treatment (a reflected rule could be abused by crafted downlink traffic).
\end{solutionbox}

\begin{solutionbox}{\ref{exr:pdb}}
Remote UPF: $600\times5\,\mu\mathrm{s} = 3$\,ms + 2\,ms = 5\,ms. 5QI~3: $50 - 5 = 45$\,ms (remote) or 49.5\,ms (local); 5QI~82: $10 - 5 = 5$\,ms (remote) or 9.5\,ms
(local). 5QI~82 loses half its access budget with the remote UPF -- it needs an edge UPF; 5QI~3 does not.
\end{solutionbox}

\begin{solutionbox}{\ref{exr:harq}}
$p = 0.1$: $n = 6$, about 6\,ms; $p = 0.001$: $n = 2$, about 2\,ms. A conservative MCS (low BLER) costs spectral efficiency but saves retransmissions and delay -- the
characteristic URLLC trade-off of capacity for latency and reliability.
\end{solutionbox}

\begin{solutionbox}{\ref{exr:adm}}
Requirement: $T_2 \le 2T_{\mathrm{FIFO}} \Leftrightarrow 1/(1-\rho_1) \le 2 \Leftrightarrow \rho_1 \le 0.5$: at most 100 calls (0.5\,\% each). In practice operators limit GBR load
much lower (e.g.\ 20--30\,\%) to keep non-GBR delay close to FIFO and to keep margin for channel variations that change the cost per call.
\end{solutionbox}

\begin{solutionbox}{\ref{exr:e2e}}
5QI: a GBR video-type value (e.g.\ 3 for real-time gaming characteristics) for the game stream, with GFBR set to the needed rate or a high-priority non-GBR 5QI with
sufficient Session-AMBR; 5QI~5/default for control. UPF at a regional edge so that the core part of the RTT is a few milliseconds; possibly an eMBB/gaming slice with
RAN resource shares in busy cells; transport mapping of the flow to a priority class. Enforcement: UPF (classification, MFBR, AMBR), gNB (GFBR, scheduling priority,
PDB), UE (uplink rules); the PCF authorises it based on the gaming provider's AF request.
\end{solutionbox}

% ---------------------------------------------------------------------
\section{Further reading}
The 5G QoS model and the standardised 5QI table are in TS~23.501 \cite{ts23501} (clause 5.7), policy and charging control in TS~23.503 \cite{ts23503}, QoS
procedures in TS~23.502 \cite{ts23502}, SDAP and RAN aspects in TS~38.300 \cite{ts38300}. The LTE counterpart is TS~23.203 \cite{ts23203}. Kleinrock
\cite{kleinrock1975queueing} derives priority queueing results; Dahlman et al.\ \cite{dahlman2020nr} discuss URLLC techniques in the NR radio.

\printbibliography[heading=subbibliography,title={\ACNstrRefs}]
\end{document}
